\documentclass[12pt]{article}
\usepackage[utf8]{inputenc}
\usepackage[T1]{fontenc}
\usepackage[margin=1in]{geometry}
\usepackage{setspace}
\usepackage{xcolor}
\usepackage[colorlinks=true, urlcolor=blue!60!black, linkcolor=blue!60!black]{hyperref}
\setlength{\parindent}{2em}
\setlength{\parskip}{1.25em}
\renewcommand{\baselinestretch}{1.0}
\begin{document}
\begin{center}
{\LARGE \textit{United States v. Lopez}, 514 U.S. 549 (1995)\par}
\vspace{0.8em}
{\large Case Brief --- Four-Page Edition\par}
\end{center}
\vspace{1em}

\section{Memory Jogger}

For the first time in nearly six decades, the Court struck down a federal statute as exceeding the Commerce Clause, holding that possessing a gun in a local school zone is neither economic activity nor tied by a jurisdictional element to interstate commerce (\href{https://www.courtlistener.com/opinion/117927/united-states-v-lopez/}{514 U.S. 549}).

\section{Detailed Case Facts}

On March 10, 1992, Alfonso Lopez, Jr., a twelfth-grader at Edison High School in San Antonio, came to school with a concealed .38-caliber handgun and five bullets. Confronted on an anonymous tip, he admitted carrying the gun and was charged under Texas law. The next day state charges were dropped after federal agents charged him under the Gun-Free School Zones Act of 1990, which made it a federal crime ``for any individual knowingly to possess a firearm at a place that the individual knows, or has reasonable cause to believe, is a school zone''---within 1,000 feet of a school. The Act as enacted contained no congressional findings linking school-zone gun possession to interstate commerce and required no proof that the gun had moved in commerce.

\section{Procedural History}

The District Court denied Lopez's motion to dismiss, reasoning that the ``business'' of schools affects interstate commerce, and after a bench trial sentenced him to six months' imprisonment and two years' supervised release. The Fifth Circuit reversed, holding that the Act exceeded the commerce power. The Supreme Court granted certiorari and affirmed on April 26, 1995.

\section{Judicial Votes}

Affirmed, \textbf{5--4}.
\begin{itemize}
\item \textbf{Majority:} Chief Justice Rehnquist, joined by Justices O'Connor, Scalia, Kennedy, and Thomas.
\item \textbf{Concurrences:} Justice Kennedy, joined by Justice O'Connor; Justice Thomas.
\item \textbf{Dissents:} Justice Stevens; Justice Souter; Justice Breyer, joined by Justices Stevens, Souter, and Ginsburg.
\end{itemize}

\section{Holding}

The Act exceeds Congress's commerce power. Congress may regulate (1) the channels of interstate commerce; (2) its instrumentalities, and persons or things in it; and (3) intrastate activities that ``substantially affect'' interstate commerce. A criminal statute regulating local, noneconomic conduct that is not part of a larger economic regulation, has no jurisdictional element, and depends on an attenuated causal chain fails the third test. Remarks about congressional findings and hypothetical federal regulation of family law or curricula are dicta, and the Court did not decide whether a jurisdictional element alone would suffice; Congress added one in 1996.

\section{Analysis of Opinions}

\textbf{Majority (Rehnquist).} Starting from ``first principles'' of enumerated powers, Rehnquist traced the doctrine from \textit{Gibbons} through \textit{Jones \& Laughlin}, \textit{Darby}, and \textit{Wickard}, finding outer limits even in modern cases. Section 922(q) failed category three: it regulated no economic enterprise, unlike \textit{Wickard}; it lacked a jurisdictional element; and without findings the Court could not see an effect ``not visible to the naked eye.'' The Government's ``costs of crime'' and ``national productivity'' theories had no limiting principle and would erase the line between ``what is truly national and what is truly local.''

\textbf{Kennedy (with O'Connor).} The holding is ``necessary though limited'': \textit{stare decisis} protects regulation of a national market, but federalism secures ``two distinct and discernable lines of political accountability,'' and education, where States serve as ``laboratories for experimentation,'' is a traditional state concern.

\textbf{Thomas.} The ``substantial effects'' test departs from the original meaning of ``commerce,'' which meant buying, selling, bartering, and transporting, and makes much of Article I, \S~8 surplusage.

\textbf{Stevens and Souter.} Stevens reasoned that guns are articles of commerce whose possession Congress may bar in particular markets. Souter warned of a return to \textit{Lochner}-era formalism and urged deferential rational-basis review.

\textbf{Breyer (with Stevens, Souter, and Ginsburg).} Considering aggregate effects and deferring to Congress, Breyer found an ample rational basis linking school gun violence to education, workforce skills, and the economy, and warned that the economic/noneconomic line was unworkable and unsettled over 100 statutes.

\section{Examples: Future Applications}

\begin{enumerate}
\item \textbf{Same-Side (Invalid): Brass Knuckles Near Parks.} A standalone ban on possessing brass knuckles near municipal parks, with no jurisdictional element or market scheme, regulates noneconomic possession and requires piling inference upon inference.
\item \textbf{Same-Side (Invalid): Federal Truancy Crime.} Punishing high-school truancy to boost national productivity regulates noneconomic conduct in education, using a rationale \textit{Lopez} rejected.
\item \textbf{Opposite-Side (Valid): Home Marijuana Cultivation.} A ban on local cultivation within a comprehensive scheme regulating the interstate drug market is ``an essential part of a larger regulation of economic activity,'' as \textit{Gonzales v. Raich} held.
\item \textbf{Opposite-Side (Valid): Channels and Instrumentalities.} Banning interstate shipment of unserialized firearms or sabotage of interstate rail lines falls within categories one and two.
\item \textbf{Fence-Sitter: Reenacted Act With a Jurisdictional Hook.} The 1996 version requires a firearm that ``has moved in or that otherwise affects interstate or foreign commerce''; under \textit{Scarborough} prior movement suffices, but critics call the hook a formalistic loophole.
\end{enumerate}

\section{Critique}

\textbf{Originalist views.} Randy Barnett defends \textit{Lopez} as a restoration of enumerated powers and argues that original meaning supports Thomas's narrower reading of ``commerce.'' Others call the majority a half-measure that kept \textit{Wickard} and an amorphous economic line, inviting \textit{Raich}.

\textbf{Progressive views.} Erwin Chemerinsky and others, echoing Souter and Breyer, argue that \textit{Lopez} revived pre-1937 formalism without a predictable standard, doubt that judicial federalism enhances liberty, and would leave the federal-state balance to the political process.

\section{Key Quotations}

\begin{quote}
``To uphold the Government's contentions here, we would have to pile inference upon inference in a manner that would bid fair to convert congressional authority under the Commerce Clause to a general police power of the sort retained by the States.''
\end{quote}
\noindent --- Chief Justice Rehnquist, Opinion of the Court.

\begin{quote}
``Section 922(q) is a criminal statute that by its terms has nothing to do with `commerce' or any sort of economic enterprise, however broadly one might define those terms. Section 922(q) is not an essential part of a larger regulation of economic activity, in which the regulatory scheme could be undercut unless the intrastate activity were regulated.''
\end{quote}
\noindent --- Chief Justice Rehnquist, Opinion of the Court.

\begin{quote}
``Were the Federal Government to take over the regulation of entire areas of traditional state concern, areas having nothing to do with the regulation of commercial activities, the boundaries between the spheres of federal and state authority would blur and political responsibility would become illusory.''
\end{quote}
\noindent --- Justice Kennedy, concurring.

\end{document}
